\begin{table}[!htbp]\centering\scriptsize\setlength{\tabcolsep}{3pt}
\caption{Franchise d'activité contre taux d'entrée : le meilleur barème pour chaque combinaison}\label{tab:ircsg_design}
\resizebox{\linewidth}{!}{%
\begin{tabular}{@{}rrrrrrrl@{}}
\toprule
Franchise & Taux & Gain & dont & Retraités & \multicolumn{2}{c}{Premier salaire (\%)} & Barème \\ \cmidrule(lr){6-7}
(€) & d'entrée & Md€ & participation & Md€ & mi-temps & Smic & \\
\midrule
0 & 0~\% & 3,26 & 0,88 & +6,9 & 4,6 & 9,2 & 22~\% à 20~k€, 22,5~\% à 30~k€, 50~\% à 212~k€ \\
0 & 2~\% & 3,49 & 0,78 & +6,7 & 5,7 & 9,7 & 21~\% à 20~k€, 21,2~\% à 30~k€, 50~\% à 194~k€ \\
0 & 4~\% & 3,69 & 0,82 & +6,6 & 6,7 & 9,8 & 19~\% à 20~k€, 24,4~\% à 45~k€, 50~\% à 201~k€ \\
0 & 6~\% & 3,89 & 0,69 & +6,4 & 7,8 & 10,3 & 18~\% à 20~k€, 21,7~\% à 37,5~k€, 50~\% à 184~k€ \\
0 & 8~\% & 4,07 & 0,76 & +6,2 & 8,8 & 10,4 & 16~\% à 20~k€, 24,5~\% à 45~k€, 50~\% à 194~k€ \\
0 & 10~\% & 4,20 & 0,96 & +6,2 & 9,5 & 10,1 & 13~\% à 20~k€, 25,2~\% à 37,5~k€, 50~\% à 219~k€ \\
2\,000 & 0~\% & 3,26 & 1,22 & +10,6 & 3,1 & 8,0 & 24~\% à 20~k€, 28,8~\% à 90~k€, 50~\% à 176~k€ \\
2\,000 & 2~\% & 3,51 & 1,28 & +10,3 & 4,0 & 8,3 & 22~\% à 20~k€, 26,3~\% à 45~k€, 50~\% à 231~k€ \\
2\,000 & 4~\% & 3,74 & 1,16 & +10,1 & 5,0 & 8,9 & 21~\% à 20~k€, 32,3~\% à 90~k€, 50~\% à 212~k€ \\
2\,000 & 6~\% & 3,96 & 1,20 & +9,9 & 5,9 & 9,1 & 19~\% à 20~k€, 24,6~\% à 37,5~k€, 50~\% à 212~k€ \\
2\,000 & 8~\% & 4,17 & 1,23 & +9,8 & 6,8 & 9,4 & 17~\% à 20~k€, 25,3~\% à 37,5~k€, 50~\% à 219~k€ \\
2\,000 & 10~\% & 4,36 & 1,26 & +9,6 & 7,6 & 9,6 & 15~\% à 20~k€, 26,0~\% à 37,5~k€, 50~\% à 227~k€ \\
2\,000 & 12~\% & 4,56 & 1,13 & +9,3 & 8,6 & 10,2 & 14~\% à 20~k€, 24,7~\% à 37,5~k€, 50~\% à 205~k€ \\
4\,000 & 0~\% & 3,30 & 1,62 & +14,4 & 1,8 & 6,8 & 26~\% à 20~k€, 28,8~\% à 90~k€, 50~\% à 182~k€ \\
4\,000 & 2~\% & 3,52 & 1,50 & +14,2 & 2,6 & 7,4 & 25~\% à 20~k€, 27,4~\% à 90~k€, 50~\% à 156~k€ \\
4\,000 & 4~\% & 3,78 & 1,54 & +13,9 & 3,4 & 7,7 & 23~\% à 20~k€, 32,0~\% à 90~k€, 50~\% à 212~k€ \\
4\,000 & 6~\% & 4,01 & 1,67 & +13,8 & 4,0 & 7,8 & 20~\% à 20~k€, 25,9~\% à 30~k€, 50~\% à 247~k€ \\
4\,000 & 8~\% & 4,23 & 1,57 & +13,4 & 4,8 & 8,4 & 19~\% à 20~k€, 26,4~\% à 37,5~k€, 50~\% à 231~k€ \\
4\,000 & 10~\% & 4,44 & 1,58 & +13,2 & 5,6 & 8,8 & 17~\% à 20~k€, 27,0~\% à 37,5~k€, 50~\% à 239~k€ \\
4\,000 & 12~\% & 4,64 & 1,45 & +12,9 & 6,4 & 9,4 & 16~\% à 20~k€, 25,7~\% à 37,5~k€, 50~\% à 215~k€ \\
4\,000 & 14~\% & 4,83 & 1,45 & +12,7 & 7,2 & 9,8 & 14~\% à 20~k€, 26,4~\% à 37,5~k€, 50~\% à 223~k€ \\
6\,000 & 6~\% & 4,05 & 1,93 & +17,6 & 2,4 & 6,9 & 23~\% à 20~k€, 25,2~\% à 25~k€, 50~\% à 239~k€ \\
6\,000 & 8~\% & 4,28 & 1,92 & +17,3 & 2,9 & 7,4 & 21~\% à 20~k€, 26,1~\% à 30~k€, 50~\% à 239~k€ \\
6\,000 & 10~\% & 4,50 & 1,91 & +17,1 & 3,5 & 7,8 & 19~\% à 20~k€, 26,2~\% à 30~k€, 50~\% à 239~k€ \\
6\,000 & 12~\% & 4,71 & 2,00 & +17,0 & 4,0 & 8,0 & 16~\% à 20~k€, 27,8~\% à 30~k€, 50~\% à 274~k€ \\
6\,000 & 14~\% & 4,89 & 1,87 & +16,6 & 4,5 & 8,7 & 15~\% à 20~k€, 26,4~\% à 30~k€, 50~\% à 239~k€ \\
7\,000 & 6~\% & 4,08 & 2,03 & +19,6 & 1,6 & 6,5 & 25~\% à 20~k€, 31,3~\% à 90~k€, 50~\% à 200~k€ \\
7\,000 & 8~\% & 4,28 & 2,01 & +19,3 & 2,0 & 6,9 & 23~\% à 20~k€, 25,5~\% à 30~k€, 50~\% à 227~k€ \\
7\,000 & 10~\% & 4,53 & 2,09 & +19,1 & 2,4 & 7,2 & 20~\% à 20~k€, 27,1~\% à 30~k€, 50~\% à 256~k€ \\
7\,000 & 12~\% & 4,73 & 2,07 & +18,9 & 2,8 & 7,7 & 18~\% à 20~k€, 27,2~\% à 30~k€, 50~\% à 256~k€ \\
7\,000 & 14~\% & 4,89 & 2,04 & +18,6 & 3,2 & 8,1 & 16~\% à 20~k€, 27,2~\% à 30~k€, 50~\% à 247~k€ \\
7\,000 & 16~\% & 5,02 & 1,80 & +18,2 & 3,7 & 8,9 & 16~\% à 20~k€, 24,3~\% à 30~k€, 50~\% à 201~k€ \\
8\,000 & 8~\% & 4,31 & 2,21 & +21,4 & 1,2 & 6,3 & 24~\% à 20~k€, 25,8~\% à 25~k€, 50~\% à 239~k€ \\
8\,000 & 10~\% & 4,51 & 2,19 & +21,1 & 1,4 & 6,8 & 22~\% à 20~k€, 26,5~\% à 30~k€, 50~\% à 239~k€ \\
8\,000 & 12~\% & 4,72 & 2,23 & +21,0 & 1,6 & 7,1 & 19~\% à 20~k€, 26,4~\% à 25~k€, 50~\% à 247~k€ \\
8\,000 & 14~\% & 4,91 & 2,27 & +20,8 & 1,9 & 7,4 & 16~\% à 20~k€, 27,4~\% à 25~k€, 50~\% à 265~k€ \\
8\,000 & 16~\% & 5,04 & 2,06 & +20,3 & 2,2 & 8,2 & 16~\% à 20~k€, 24,7~\% à 25~k€, 50~\% à 215~k€ \\
\bottomrule
\end{tabular}}
\par\vspace{3pt}\parbox{\linewidth}{\raggedright\scriptsize\color{encre}
Pour chaque franchise d'activité (par adulte) et chaque taux d'entrée (taux sur le premier euro imposé après la franchise), le barème le plus efficace parmi ceux qui respectent la recette, les déciles à $\pm$500~€ et la neutralité du 1~\%. Sont aussi fixées : l'exonération des petites pensions, la franchise de 6\,000~€ par enfant, la déduction de l'épargne retraite. Premier salaire : taux de prélèvement IR-CSG d'une personne seule qui prend un emploi au Smic, à mi-temps ou à temps plein (aujourd'hui : 9,5~\%). Taux d'entrée pairs seulement ; 10332 barèmes essayés, 1047 admissibles. La proposition retient la combinaison qui maximise la participation (franchise de 8\,000~€, taux d'entrée de 14~\%).}
\end{table}
